% ================================================================
% PROCESSI PRIMARI
% ================================================================

\section{Processi Primari}
\label{sec:processi-primari}

In conformità allo standard ISO/IEC~12207:1995, i processi primari
comprendono le attività direttamente coinvolte nella realizzazione, nella fornitura
e nell'evoluzione del prodotto software.

Nel contesto del progetto \textit{Second Brain\textsubscript{G}}, e in accordo con il tailoring
descritto nella Sezione~\ref{sec:tailoring}, il gruppo \textit{Synergo Unit} adotta
i seguenti processi primari:

\begin{itemize}
    \item \textbf{Fornitura}: disciplina il rapporto tra il gruppo, il committente e
    la proponente, regolando candidatura, pianificazione, produzione degli
    artefatti, consegna e comunicazioni ufficiali;

    \item \textbf{Sviluppo}: disciplina le attività tecniche necessarie alla
    definizione e alla progressiva realizzazione del prodotto software, con
    particolare riferimento, in fase RTB, all'analisi dei requisiti,
    all'analisi dello stato dell'arte, alla progettazione preliminare e all'implementazione dimostrativa del Proof of Concept.
\end{itemize}

Le attività di codifica estesa, testing automatico, integrazione continua e
rilascio software sono state normate progressivamente nella fase PB, in
seguito alla stabilizzazione dell'architettura e all'avvio dello sviluppo
del prodotto; sono descritte in dettaglio nelle sottosezioni seguenti e nel
Processo di Verifica (Sezione~\ref{sec:verifica}).

% ================================================================
% FORNITURA
% ================================================================

\subsection{Processo di Fornitura}
\label{sec:fornitura}

\subsubsection{Descrizione}

Il processo di fornitura regola le attività mediante cui il gruppo
\textit{Synergo Unit} opera come fornitore nei confronti del committente e della
proponente.

Nel contesto del progetto didattico:

\begin{itemize}
    \item il \textbf{docente} ricopre il ruolo\textsubscript{G} di committente rispetto alle forniture,
    alle scadenze, ai vincoli regolamentari e alla valutazione delle baseline;

    \item la \textbf{proponente}, Zucchetti S.p.A., ricopre il ruolo
    di cliente rispetto alle esigenze di prodotto e di mentore rispetto alle
    scelte tecniche e progettuali;

    \item il gruppo \textit{Synergo Unit} ricopre il ruolo di fornitore, assumendosi
    la responsabilità di produrre, verificare e consegnare gli artefatti richiesti
    dalle revisioni di avanzamento.
\end{itemize}

Il processo comprende:

\begin{itemize}
    \item la valutazione dei capitolati\textsubscript{G};
    \item la candidatura al capitolato scelto;
    \item la formalizzazione degli impegni economici e temporali;
    \item la produzione della documentazione richiesta per ciascuna baseline;
    \item la gestione delle comunicazioni ufficiali;
    \item la rendicontazione dell'avanzamento;
    \item la consegna della documentazione tramite repository\textsubscript{G} pubblico e sito web.
\end{itemize}

% ------------------------------------------------

\subsubsection{Norme e Strumenti del Processo di Fornitura}

\paragraph{Strumenti a Supporto}

Il processo di fornitura è supportato dai seguenti strumenti:

\begin{itemize}
      \item \textbf{Email istituzionale del gruppo}: l'indirizzo
            \texttt{synergounit.20@gmail.com} viene utilizzato per le comunicazioni ufficiali
            con committente e proponente, garantendo tracciabilità e continuità nelle
            interazioni esterne;

      \item \textbf{Repository GitHub}: utilizzato come punto di raccolta, versionamento
        e pubblicazione degli artefatti documentali e del codice del Proof of Concept;

      \item \textbf{GitHub Pages\textsubscript{G}}: utilizzato per rendere consultabile la documentazione
      prodotta;

      \item \textbf{GitHub Projects}: utilizzato per il tracciamento delle attività,
      delle milestone\textsubscript{G} RTB e PB e dello stato di avanzamento;

      \item \textbf{Google Sheets}: utilizzato per la raccolta e la consuntivazione
      delle ore previste ed effettive;

      \item \textbf{Google Forms}: utilizzato per la creazione standardizzata delle
      attività e il successivo tracciamento su GitHub;

      \item \textbf{Canva}: utilizzato per la predisposizione del materiale di supporto
      ai Diari di Bordo.
\end{itemize}

Le modalità operative di utilizzo degli strumenti sono dettagliate nei processi
di supporto e nei processi organizzativi.

% ------------------------------------------------

\paragraph{Documentazione prodotta in fase CC}

Durante la fase di candidatura (CC), il gruppo ha prodotto e
approvato la documentazione necessaria alla presentazione della propria proposta
per il capitolato C6 -- \textit{Second Brain}.

\bigskip
\begin{xltabular}{\textwidth}{>{\hsize=.32\hsize\raggedright\arraybackslash}X
                                >{\hsize=.68\hsize\raggedright\arraybackslash}X}
    \toprule
    \textbf{Documento} & \textbf{Descrizione} \\
    \midrule
    \endfirsthead

    \toprule
    \textbf{Documento} & \textbf{Descrizione} \\
    \midrule
    \endhead

    \midrule
    \multicolumn{2}{r}{\textit{Continua nella pagina successiva}}\\
    \endfoot

    \bottomrule
    \endlastfoot

    \doclink{https://github.com/Synergo-Unit-UniPD/Second-Brain/blob/main/docs/CC/doc_gestionali/doc_esterni/lettera_presentazione/lettera_presentazione.pdf}
    {Lettera di Presentazione}\textsubscript{G}
    & Documento con cui il gruppo comunica formalmente la candidatura al capitolato
    scelto. Contiene l'indicazione del capitolato, della proponente, della scadenza
    prevista, del budget\textsubscript{G} preventivato e degli impegni assunti dal gruppo. \\

    \doclink{https://github.com/Synergo-Unit-UniPD/Second-Brain/blob/main/docs/CC/doc_gestionali/doc_esterni/valutazione_capitolati/valutazione_capitolati.pdf}
    {Valutazione dei Capitolati}
    & Documento esterno che raccoglie il confronto tra i capitolati analizzati,
    evidenziando per ciascuno aspetti positivi, criticità, rischi e motivazioni
    della scelta finale. \\

    \doclink{https://github.com/Synergo-Unit-UniPD/Second-Brain/blob/main/docs/CC/doc_gestionali/doc_esterni/preventivo_costi_assunzione_impegni/preventivo_costi_e_assunzione_impegni.pdf}
    {Preventivo Costi e Assunzione Impegni}
    & Documento esterno che formalizza il preventivo economico, la distribuzione
    delle ore produttive, la ripartizione del lavoro per ruolo e gli impegni assunti
    dal gruppo verso il committente. \\

    Norme di Progetto
    & Documento interno, approvato in versione 1.x.x per la baseline CC, contenente
    le norme operative iniziali del gruppo. \\

    Glossario
    & Documento interno, approvato in versione 1.x.x per la baseline CC, contenente
    i termini tecnici e ambigui utilizzati nella documentazione iniziale. \\

    Verbali interni
    & Documenti interni che tracciano decisioni organizzative, assegnazioni di attività
    e discussioni svolte dal gruppo. \\

    Verbali esterni
    & Documenti esterni che formalizzano gli incontri con le aziende proponenti e i
    chiarimenti ottenuti durante la fase di candidatura. \\

\end{xltabular}

\bigskip

La versione 1.x.x delle Norme di Progetto e del Glossario è relativa alla sola
fase di candidatura. Con l'ingresso nella fase RTB, tali documenti vengono
aggiornati alla versione 2.x.x per recepire i nuovi processi attivati.

% ------------------------------------------------

\paragraph{Documentazione prodotta o prevista in fase RTB}

Durante la fase RTB, il gruppo produce e aggiorna la documentazione necessaria
alla Requirements \& Technology Baseline\textsubscript{G}.

\bigskip

\begin{xltabular}{\textwidth}{
    >{\hsize=.23\hsize\raggedright\arraybackslash}X
    >{\hsize=.18\hsize\raggedright\arraybackslash}X
    >{\hsize=.12\hsize\raggedright\arraybackslash}X
    >{\hsize=.47\hsize\raggedright\arraybackslash}X
}
    \toprule
    \textbf{Documento} & \textbf{Redattore\textsubscript{G}} & \textbf{Uso} & \textbf{Descrizione} \\
    \midrule
    \endfirsthead

    \toprule
    \textbf{Documento} & \textbf{Redattore} & \textbf{Uso} & \textbf{Descrizione} \\
    \midrule
    \endhead

    \midrule
    \multicolumn{4}{r}{\textit{Continua nella pagina successiva}}\\
    \endfoot

    \bottomrule
    \endlastfoot

    Lettera di Presentazione RTB
    & Responsabile\textsubscript{G}
    & Esterno
    & Documento redatto prima della candidatura alla RTB. Formalizza la richiesta di
    partecipazione alla revisione\textsubscript{G} e rimanda alla documentazione disponibile nel
    repository pubblico. \\

    \doclink{https://synergo-unit-unipd.github.io/Second-Brain/docs/PB/doc_gestionali/doc_esterni/piano_di_progetto/piano_di_progetto.pdf}
    {Piano di Progetto}
    & Amministratore\textsubscript{G}
    & Esterno
    & Documento che descrive pianificazione, consuntivazione, gestione dei rischi,
    rotazione dei ruoli, avanzamento degli sprint\textsubscript{G} e controllo dei costi. \\

    \doclink{https://synergo-unit-unipd.github.io/Second-Brain/docs/PB/doc_gestionali/doc_esterni/piano_di_qualifica/piano_di_qualifica.pdf}
    {Piano di Qualifica}
    & Amministratore
    & Esterno
    & Documento che definisce strategie, metriche e attività di verifica e validazione
    adottate dal gruppo per garantire la qualità di processo e di prodotto. \\

    \doclink{https://synergo-unit-unipd.github.io/Second-Brain/docs/PB/doc_tecnici/doc_esterni/analisi_dei_requisiti/analisi_dei_requisiti.pdf}
    {Analisi dei Requisiti}
    & Analista\textsubscript{G}
    & Esterno
    & Documento tecnico che formalizza casi d'uso\textsubscript{G}, requisiti, fonti e tracciabilità
    tra esigenze espresse, funzionalità attese e vincoli del prodotto. \\

    Norme di Progetto
    & Amministratore
    & Interno
    & Documento che definisce il way of working del gruppo. In RTB viene
    aggiornato alla versione 2.x.x per includere i processi attivi dopo la candidatura. \\

    \doclink{https://synergo-unit-unipd.github.io/Second-Brain/docs/PB/doc_gestionali/doc_interni/glossario/glossario.pdf}
    {Glossario}
    & Amministratore
    & Interno
    & Documento contenente termini tecnici, acronimi ed espressioni ambigue. Ogni
    autore propone i termini emersi durante la stesura; l'Amministratore ne cura
    normalizzazione e coerenza. \\

    \doclink{https://synergo-unit-unipd.github.io/Second-Brain/docs/PB/doc_gestionali/doc_interni/analisi_stato_arte/analisi_stato_arte.pdf}
    {Analisi dello Stato dell'Arte}
    & Analista / Progettista\textsubscript{G}
    & Interno
    & Documento che raccoglie la valutazione comparativa delle tecnologie e la scelta
    dello stack candidato. L'Analista cura la valutazione; il Progettista seleziona
    le tecnologie sulla base delle esigenze architetturali. \\

    Verbali interni
    & Responsabile
    & Interno
    & Documenti che tracciano riunioni, decisioni operative, criticità, retrospettive
    e pianificazione interna. \\

    Verbali esterni
    & Amministratore
    & Esterno
    & Documenti che formalizzano incontri con la proponente e chiarimenti rilevanti
    ai fini dell'Analisi dei Requisiti e delle scelte tecnologiche. \\

    Diari di Bordo
    & Responsabile
    & Esterno
    & Slide e materiale gestionale utilizzati per il confronto periodico con il
    committente, con evidenza di avanzamento, criticità e obiettivi successivi. \\

\end{xltabular}

\bigskip

I destinatari dei documenti esterni sono:

\begin{itemize}
    \item il gruppo \textit{Synergo Unit};
    \item Zucchetti S.p.A.;
    \item il Prof.\ Tullio Vardanega;
    \item il Prof.\ Riccardo Cardin.
\end{itemize}

I Diari di Bordo devono riportare almeno:
\begin{itemize}
    \item risultati raggiunti nel periodo di riferimento;
    \item obiettivi previsti per il periodo successivo;
    \item criticità emerse;
    \item eventuali scostamenti rispetto alla pianificazione;
    \item decisioni o richieste di chiarimento da sottoporre al committente.
\end{itemize}

Tale materiale ha funzione di supporto alla comunicazione periodica con il
committente e non sostituisce la documentazione ufficiale di progetto.

Il Proof of Concept non prevede, nella fase RTB, un documento separato dedicato:
le sue finalità, il perimetro e gli esiti verranno descritti all'interno del materiale
di presentazione della Requirements \& Technology Baseline e, ove necessario, nei Diari
di Bordo.

Gli Architecture Decision Records (ADR) non sono ancora considerati
artefatti attivi della fase RTB; la loro introduzione è pianificata per la fase PB,
in seguito alla validazione preliminare delle scelte tecnologiche e architetturali.

% ------------------------------------------------

\paragraph{Documentazione prodotta in fase PB}

Durante la fase PB, il gruppo aggiorna la documentazione già prodotta in RTB e ne
produce di nuova, necessaria alla Product Baseline\textsubscript{G} e coerente con il
completamento del Minimum Viable Product.

\bigskip

\begin{xltabular}{\textwidth}{
    >{\hsize=.23\hsize\raggedright\arraybackslash}X
    >{\hsize=.18\hsize\raggedright\arraybackslash}X
    >{\hsize=.12\hsize\raggedright\arraybackslash}X
    >{\hsize=.47\hsize\raggedright\arraybackslash}X
}
    \toprule
    \textbf{Documento} & \textbf{Redattore} & \textbf{Uso} & \textbf{Descrizione} \\
    \midrule
    \endfirsthead

    \toprule
    \textbf{Documento} & \textbf{Redattore} & \textbf{Uso} & \textbf{Descrizione} \\
    \midrule
    \endhead

    \midrule
    \multicolumn{4}{r}{\textit{Continua nella pagina successiva}}\\
    \endfoot

    \bottomrule
    \endlastfoot

    \doclink{https://synergo-unit-unipd.github.io/Second-Brain/docs/PB/doc_gestionali/doc_esterni/piano_di_progetto/piano_di_progetto.pdf}
    {Piano di Progetto}
    & Amministratore
    & Esterno
    & Aggiornato con il consuntivo RTB, il Preventivo a Finire e la pianificazione
    delle attività fino al completamento della PB. \\

    \doclink{https://synergo-unit-unipd.github.io/Second-Brain/docs/PB/doc_gestionali/doc_esterni/piano_di_qualifica/piano_di_qualifica.pdf}
    {Piano di Qualifica}
    & Amministratore
    & Esterno
    & Aggiornato con i risultati delle verifiche e validazioni condotte sul Minimum
    Viable Product, incluse le metriche di copertura dei test. \\

    \doclink{https://synergo-unit-unipd.github.io/Second-Brain/docs/PB/doc_tecnici/doc_esterni/analisi_dei_requisiti/analisi_dei_requisiti.pdf}
    {Analisi dei Requisiti}
    & Analista
    & Esterno
    & Raffinata a seguito degli esiti del Proof of Concept e del confronto con la
    proponente durante la RTB. \\

    \doclink{https://synergo-unit-unipd.github.io/Second-Brain/docs/PB/doc_tecnici/doc_esterni/specifica_tecnica/specifica_tecnica.pdf}
    {Specifica Tecnica}
    & Progettista
    & Esterno
    & Documento tecnico introdotto in fase PB: descrive architettura, decomposizione
    in componenti, design pattern adottati, flusso dei dati e tracciamento
    componenti-requisiti del prodotto finale. \\

    \doclink{https://synergo-unit-unipd.github.io/Second-Brain/docs/PB/doc_tecnici/doc_esterni/manuale_utente/manuale_utente.pdf}
    {Manuale Utente}
    & Progettista
    & Esterno
    & Documento introdotto in fase PB: descrive in linguaggio non tecnico
    installazione, configurazione e utilizzo di ciascuna funzionalità del prodotto,
    inclusa la risoluzione dei problemi più comuni. \\

    Norme di Progetto
    & Amministratore
    & Interno
    & Aggiornato progressivamente durante la PB (dalla versione 2.x.x) per
    recepire i processi di sviluppo, verifica e rilascio formalizzati in questa
    fase (Sezione~\ref{sec:tailoring}). \\

    \doclink{https://synergo-unit-unipd.github.io/Second-Brain/docs/PB/doc_gestionali/doc_interni/glossario/glossario.pdf}
    {Glossario}
    & Amministratore
    & Interno
    & Ampliato con la terminologia tecnica emersa durante l'implementazione del
    Minimum Viable Product. \\

    Verbali interni
    & Responsabile
    & Interno
    & Documenti che tracciano riunioni, decisioni operative, criticità e
    retrospettive di sprint svolte durante la fase PB. \\

    Verbali esterni
    & Amministratore
    & Esterno
    & Documenti che formalizzano gli incontri con la proponente e i chiarimenti
    rilevanti ai fini della Specifica Tecnica e del Manuale Utente. \\

    Diari di Bordo
    & Responsabile
    & Esterno
    & Slide e materiale gestionale utilizzati per il confronto periodico con il
    committente durante la fase PB, con evidenza di avanzamento, criticità e
    obiettivi successivi. \\

\end{xltabular}

% ------------------------------------------------

\subsubsection{Attività del Processo di Fornitura}

\paragraph{Inizializzazione}

\begin{itemize}
    \item \textbf{Ruoli}: Responsabile, Amministratore.
    \item \textbf{Input}: capitolati disponibili, regolamento del progetto didattico,
    indicazioni dei docenti.
    \item \textbf{Output}: decisione di partecipazione alla selezione del capitolato.
\end{itemize}

L'attività consiste nell'analisi preliminare delle proposte disponibili, nella
valutazione della compatibilità tra competenze del gruppo e richieste del capitolato
e nella decisione di procedere con la candidatura.

\paragraph{Risposta}

\begin{itemize}
    \item \textbf{Ruoli}: Responsabile.
    \item \textbf{Input}: capitolato selezionato, valutazione dei capitolati, preventivo.
    \item \textbf{Output}: Lettera di Presentazione.
\end{itemize}

Il gruppo formalizza la propria candidatura mediante una Lettera di Presentazione,
nella quale dichiara il capitolato scelto, la motivazione della scelta, la scadenza
prevista, il budget preventivato e gli impegni assunti.

\paragraph{Consolidamento degli impegni}

\begin{itemize}
    \item \textbf{Ruoli}: Responsabile, Amministratore.
    \item \textbf{Input}: candidatura accettata, Preventivo Costi e Assunzione Impegni.
    \item \textbf{Output}: impegni economici e temporali consolidati.
\end{itemize}

A seguito dell'aggiudicazione, il gruppo considera vincolanti gli impegni dichiarati
in fase di candidatura.

Gli impegni assunti prevedono:
\begin{itemize}
    \item completamento del progetto entro il 2026-08-28;
    \item rispetto del budget preventivato pari a 11.665,00~\euro{};
    \item mantenimento del monte ore produttivo dichiarato;
    \item rotazione dei ruoli secondo quanto pianificato nel Piano di Progetto;
    \item comunicazione costante e tracciabile con committente e proponente.
\end{itemize}

Il costo preventivato costituisce il limite massimo previsto per l'intero
svolgimento del progetto.

\paragraph{Pianificazione}

\begin{itemize}
    \item \textbf{Ruoli}: Amministratore, Responsabile.
    \item \textbf{Input}: impegni approvati, milestone di progetto, baseline prevista.
    \item \textbf{Output}: Piano di Progetto aggiornato.
\end{itemize}

L'Amministratore redige e aggiorna il Piano di Progetto, includendo pianificazione,
preventivi, consuntivi di sprint, rischi, rotazione dei ruoli e avanzamento rispetto
alle milestone.

\paragraph{Esecuzione e controllo}

\begin{itemize}
    \item \textbf{Ruoli}: tutti i ruoli attivi nello sprint.
    \item \textbf{Input}: Piano di Progetto, attività pianificate, issue\textsubscript{G} assegnate.
    \item \textbf{Output}: documentazione prodotta, consuntivi aggiornati, avanzamento
    monitorato.
\end{itemize}

Durante ogni sprint il gruppo svolge le attività pianificate, consuntiva le ore
effettive e aggiorna lo stato delle issue. Il controllo dell'avanzamento avviene
tramite GitHub Projects e fogli di consuntivazione condivisi.

\paragraph{Revisione}

\begin{itemize}
    \item \textbf{Ruoli}: Responsabile, Verificatore\textsubscript{G}, ruoli coinvolti negli artefatti.
    \item \textbf{Input}: documenti o materiali prodotti, feedback interno o esterno.
    \item \textbf{Output}: artefatti corretti, feedback integrato, eventuali nuove issue.
\end{itemize}

Il gruppo revisiona periodicamente gli artefatti prodotti tramite verifiche interne,
retrospettive di sprint, incontri con la proponente e Diari di Bordo con il committente.
Eventuali osservazioni vengono trasformate in attività tracciate.

\paragraph{Validazione esterna}

\begin{itemize}
    \item \textbf{Ruoli}: Responsabile, Analista, Progettista.
    \item \textbf{Input}: requisiti, dubbi interpretativi, scelte tecnologiche preliminari,
    materiale di supporto.
    \item \textbf{Output}: chiarimenti formalizzati, verbale esterno approvato,
      eventuali modifiche all'Analisi dei Requisiti, all'Analisi dello Stato dell'Arte
      o al perimetro del Proof of Concept.
\end{itemize}

La validazione esterna ha lo scopo di verificare con la proponente la correttezza
delle interpretazioni assunte dal gruppo, in particolare per gli aspetti riguardanti
il perimetro funzionale, l'interazione con LLM, la gestione delle note Markdown,
il distant writing\textsubscript{G} e i requisiti prestazionali.

Ogni chiarimento rilevante deve essere riportato in un verbale esterno e, se
necessario, tradotto in modifiche tracciate alla documentazione.

I verbali esterni devono essere condivisi con la proponente per confermare la
correttezza dei contenuti riportati e degli accordi emersi.

\paragraph{Consegna e completamento}

\begin{itemize}
    \item \textbf{Ruoli}: Responsabile, Amministratore.
    \item \textbf{Input}: documentazione verificata e approvata.
    \item \textbf{Output}: baseline pubblicata e candidatura alla revisione.
\end{itemize}

La consegna della documentazione avviene tramite pubblicazione nel repository
pubblico e nel sito web associato. La candidatura alla revisione viene effettuata
tramite email, riportando il puntatore alla documentazione aggiornata, senza invio
di allegati salvo diversa indicazione del committente.

% ================================================================
% SVILUPPO
% ================================================================

\subsection{Processo di Sviluppo}
\label{sec:sviluppo}

\subsubsection{Descrizione}

Il processo di sviluppo disciplina le attività tecniche che portano alla definizione
e alla progressiva realizzazione del prodotto software.

Il processo di sviluppo adottato dal gruppo è istanziato a partire dal processo
\textit{Development} definito dallo standard ISO/IEC 12207.

Esso comprende le attività riportate di seguito, che vengono attivate
progressivamente durante il ciclo di vita del progetto in funzione dello stato
di avanzamento della fornitura.

\begin{itemize}
  \item Implementazione del processo (Process implementation);
  \item Analisi dei requisiti di sistema (System requirements analysis);
  \item Progettazione dell'architettura di sistema (System architectural design);
  \item Analisi dei requisiti software (Software requirements analysis);
  \item Progettazione dell'architettura software (Software architectural design);
  \item Progettazione dettagliata del software (Software detailed design);
  \item Codifica e test del software (Software coding and testing);
  \item Integrazione del software (Software integration);
  \item Test di qualifica del software (Software qualification testing);
  \item Integrazione del sistema (System integration);
  \item Test di qualifica del sistema (System qualification testing);
  \item Installazione del software (Software installation);
  \item Supporto all'accettazione del software (Software acceptance support).
\end{itemize}

Con l'avvio dello Sprint 9, a partire dal 2026-06-02, il processo di sviluppo
include attività di codifica finalizzate alla realizzazione del Proof of
Concept.

Tali attività hanno lo scopo di validare tecnicamente le scelte effettuate,
dimostrare la fattibilità dei principali flussi funzionali individuati per la
Requirements \& Technology Baseline e fornire evidenze utili alla
successiva progettazione di dettaglio.

% ------------------------------------------------

\subsubsection{Implementazione del processo}

Il gruppo adotta un modello organizzativo iterativo ispirato a Scrum,
basato su sprint settimanali con cadenza da martedì a martedì.

Il workflow prevede:

\begin{itemize}
    \item sprint review\textsubscript{G} dello sprint precedente;
    \item sprint retrospective\textsubscript{G} dello sprint precedente;
    \item sprint planning\textsubscript{G} dello sprint corrente;
    \item daily stand-up\textsubscript{G} asincrono;
    \item riunione intermedia di allineamento;
    \item monitoraggio continuo delle attività.
\end{itemize}

A valle delle scelte iniziali di ciascuna fase, la gestione delle attività di
progetto segue ciclicamente, ad ogni periodo di avanzamento, il medesimo
ordine: \textbf{consuntivo di periodo} (azioni svolte e consumi rilevati)
$\rightarrow$ \textbf{retrospettiva} (grado di raggiungimento degli obiettivi
attesi e ragioni presunte) $\rightarrow$ \textbf{misure correttive} (con
aggiornamento dell'analisi dei rischi) $\rightarrow$ \textbf{miglioramento
della pianificazione futura} (a breve e a lungo termine) $\rightarrow$
\textbf{preventivo a finire}. È questo ciclo di controllo a governare l'avanzamento reale
del progetto: un modello incrementale, per sua natura, avanzerebbe in
un'unica direzione, aggiungendo senza mai rivedere o correggere quanto già
deciso; il gruppo, al contrario, rivede esplicitamente a ogni ciclo tanto le
stime future quanto le scelte pregresse, alla luce dei risultati e delle
criticità emerse.

La pianificazione dettagliata degli sprint, la rotazione dei ruoli e le
milestone sono riportate nel \link{https://synergo-unit-unipd.github.io/Second-Brain/docs/PB/doc_gestionali/doc_esterni/piano_di_progetto/piano_di_progetto.pdf}{Piano di Progetto}.

Le decisioni tecniche e architetturali vengono raffinate iterativamente, secondo lo stesso ciclo, sulla base di:

\begin{itemize}
    \item feedback della proponente;
    \item risultati del Proof of Concept;
    \item criticità emerse durante gli sprint;
    \item evoluzione dei requisiti;
    \item vincoli rilevati durante l'analisi dello stato dell'arte.
\end{itemize}

Tale approccio consente di evitare scelte premature e di mantenere coerenza tra
requisiti, tecnologie e architettura.

% ------------------------------------------------

\subsubsection{Analisi dei Requisiti di sistema}

\paragraph{Scopo}

L'analisi dei requisiti di sistema ha lo scopo di identificare, classificare, tracciare e mantenere
aggiornati i requisiti del progetto Second Brain.

\paragraph{Raccolta dei requisiti}

I requisiti vengono raccolti a partire da:

\begin{itemize}
    \item capitolato d'appalto C6 -- \textit{Second Brain};
    \item incontri con la proponente (verbali esterni);
    \item analisi del dominio applicativo.
\end{itemize}

\paragraph{Classificazione}

I requisiti vengono classificati secondo:

\begin{itemize}
    \item tipologia;
    \item importanza;
    \item fonte;
    \item casi d'uso associati.
\end{itemize}

\paragraph{Gestione delle modifiche}

Le modifiche ai requisiti devono:

\begin{itemize}
    \item essere discusse internamente;
    \item essere tracciate tramite issue;
    \item mantenere consistenti le matrici di tracciabilità;
    \item essere riportate nei verbali quando derivano da incontri esterni;
    \item essere validate dalla proponente quando incidono sul perimetro funzionale
    o sui vincoli del prodotto.
\end{itemize}

% ------------------------------------------------

\subsubsection{Progettazione dell'architettura di sistema}

\paragraph{Scopo}

La progettazione dell'architettura di sistema ha lo scopo di identificare le componenti hardware, software e le operazioni manuali

Il nostro progetto è prettamente software e non identifica componenti hardware.

% ------------------------------------------------

\subsubsection{Analisi dei requisiti software}

\paragraph{Scopo}

L'analisi dei requisiti software ha lo scopo di stabilire e documentare le specifiche tecniche, funzionali, di qualità e di interfaccia necessarie alla progettazione e alla realizzazione di ogni componente software del sistema



\paragraph{Nomenclatura dei requisiti}

Ogni requisito deve essere identificato tramite un codice univoco nel formato:

\begin{center}
\texttt{R[ID]-[Tipo]-[Importanza]}
\end{center}

dove:
\begin{itemize}
    \item \textbf{ID}: identificativo numerico progressivo;
    \item \textbf{Tipo}: indica la categoria del requisito:
    \begin{itemize}
        \item \texttt{F}: funzionale;
        \item \texttt{Q}: qualità;
        \item \texttt{V}: vincolo;
        \item \texttt{P}: prestazionale;
    \end{itemize}
    \item \textbf{Importanza}: indica la priorità del requisito:
    \begin{itemize}
        \item \texttt{O}: obbligatorio;
        \item \texttt{D}: desiderabile;
        \item \texttt{Op}: opzionale.
    \end{itemize}
\end{itemize}

La nomenclatura deve essere mantenuta stabile per tutta la durata della baseline.
Eventuali modifiche agli identificativi devono essere motivate, tracciate e
verificate per non compromettere la consistenza delle matrici di tracciabilità.

\paragraph{Tracciabilità}

Il gruppo mantiene matrici di tracciabilità tra:

\begin{itemize}
    \item fonti e requisiti;
    \item requisiti e casi d'uso;
    \item casi d'uso e requisiti.
\end{itemize}

La tracciabilità deve garantire:

\begin{itemize}
    \item copertura completa delle fonti;
    \item assenza di requisiti privi di motivazione;
    \item coerenza tra requisiti e casi d'uso;
    \item possibilità di risalire da ogni requisito alla fonte che lo giustifica;
    \item possibilità di risalire da ogni caso d'uso ai requisiti soddisfatti;
    \item supporto alla successiva attività di verifica e validazione.
\end{itemize}

\paragraph{Gestione dei Casi d'Uso}

I casi d'uso costituiscono il principale strumento di descrizione delle interazioni
tra utente e sistema nella fase RTB.

Essi sono organizzati nelle seguenti macroaree funzionali:

\begin{itemize}
    \item \textbf{Gestione editor e Markdown}: scrittura, formattazione\textsubscript{G} e rendering;

    \item \textbf{Interazione AI e LLM}: generazione di suggerimenti, riscrittura,
    distant writing, applicazione di stili e supporto tramite tecniche come i
    \textit{Sei Cappelli Per Pensare\textsubscript{G}};

    \item \textbf{Gestione I/O e persistenza}: apertura, salvataggio, importazione ed
    esportazione delle note.
\end{itemize}

\paragraph{Convenzioni dei casi d'uso}

I casi d'uso adottano le seguenti convenzioni:

\begin{itemize}
      \item \textbf{Denominazione}: ogni caso d'uso deve seguire il formato \texttt{UC X: Titolo}, dove \texttt{X} rappresenta l'identificativo numerico
      progressivo e \texttt{Titolo} descrive sinteticamente l'interazione modellata;

      \item \textbf{Sottocasi}: eventuali casi d'uso derivati devono mantenere una numerazione gerarchica coerente con il caso d'uso principale;
      
      \item \textbf{Ordinamento}: i diagrammi vengono mantenuti in ordine numerico e
      aggiornati in modo incrementale;

      \item \textbf{Strumenti}: i diagrammi UML\textsubscript{G} vengono realizzati e
      mantenuti tramite LucidChart/LucidSpark\textsubscript{G} o strumenti equivalenti condivisi dal
      gruppo;

      \item \textbf{Navigabilità}: i riferimenti interni tra casi d'uso devono essere
      realizzati tramite collegamenti interni al documento;

      \item \textbf{Inclusioni ed estensioni}: le relazioni tra casi d'uso devono essere
      esplicitate sia nei diagrammi sia nella descrizione testuale.
\end{itemize}

% ------------------------------------------------

\subsubsection{Progettazione dell'architettura software}

\paragraph{Scopo}

La progettazione dell'architettura software ha lo scopo di individuare una prima
struttura del sistema, coerente con i requisiti e con le tecnologie selezionate.
Tale attività trasforma i requisiti software in un'architettura
che ne descrive la struttura top-level e identifica i componenti software necessari.

In fase RTB, tale attività ha carattere esplorativo e mira a:

\begin{itemize}
    \item individuare i principali moduli del sistema;
    \item valutare le interazioni tra frontend, backend e servizi AI;
    \item giustificare la selezione dello stack tecnologico sulla base dell'Analisi dello Stato dell'Arte;
    \item evidenziare rischi tecnici;
    \item preparare il perimetro del Proof of Concept;
    \item raccogliere elementi utili per la progettazione di dettaglio in PB.
\end{itemize}

\paragraph{Integrazione dello Stato dell'Arte e selezione tecnologica}

La definizione dell'architettura si fonda sui risultati dell'Analisi dello Stato dell'Arte,
che funge da motore decisionale per la scelta dei componenti. L'architettura viene valutata secondo criteri di appropriatezza
dei metodi e degli standard di progettazione.

Il processo di selezione prevede la collaborazione tra:

\begin{itemize}
    \item Analista: cura la valutazione comparativa delle alternative tecnologiche basandosi su maturità, supporto della comunità e compatibilità con il capitolato;
    \item Progettista: seleziona le tecnologie (attualmente individuate in Vue.js, FastAPI e Docker) che meglio rispondono alle esigenze di modularità e integrazione dei servizi AI individuate nel design.
\end{itemize}

Le decisioni architetturali definitive sono state consolidate nella fase PB. Gli
ADR sono stati introdotti a valle della stabilizzazione del processo decisionale
architetturale e vengono utilizzati per documentare in modo tracciabile le scelte
rilevanti, le alternative valutate e le motivazioni della decisione finale --
ne è un esempio ADR-01, che motiva il raffinamento del comportamento
dell'intestazione Markdown a un singolo livello (R25-F-O) rispetto alla
formulazione originaria in fase RTB.

\paragraph{Stile Architetturale (Fase PB)}
A seguito del consolidamento delle decisioni nella fase PB, il sistema adotta un'architettura ibrida a due stili combinati, separati da un confine REST:
\begin{itemize}
    \item \textbf{Frontend:} adotta il pattern MVC (Model-View-Controller) con Observer in modalità \textit{pull}. Questa scelta risolve la coordinazione tra editor, anteprima e proposte AI senza accoppiare le View tra loro.
    \item \textbf{Backend:} adotta una \textit{Layered Architecture} (Presentation / Business / Persistence). Tale scelta garantisce l'indipendenza del dominio (AI, orchestrazione note) rispetto al provider LLM e al tipo di storage, limitando la complessità rispetto a un'architettura esagonale, come concordato con la proponente.
\end{itemize}

% ------------------------------------------------

\subsubsection{Progettazione dettagliata del software}

\paragraph{Scopo}

La progettazione dettagliata del software definisce il design di ogni unità software in modo che possa essere codificata. Nella fase PB, questa attività stabilisce i pattern implementativi e la gestione delle interazioni tra componenti.

\paragraph{Gestione delle Dipendenze (Dependency Injection)}

Per rispettare il principio \textit{Composition over Inheritance}, nessuna classe applicativa crea le proprie dipendenze direttamente.
\begin{itemize}
    \item \textbf{Backend:} l'iniezione delle dipendenze avviene tramite la funzione \texttt{Depends()} di FastAPI, con provider centralizzati nel package \texttt{app/di/}.
    \item \textbf{Frontend:} in assenza di un container DI nativo in Vue, l'iniezione è manuale e centralizzata nel modulo \texttt{src/bootstrap/composition.ts}, mantenendo il dominio TypeScript puro e isolato dal framework.
\end{itemize}

% ------------------------------------------------

\subsubsection{Codifica e test del software}

\paragraph{Scopo}

La codifica e test del software riguarda lo sviluppo delle unità software e delle relative procedure di test.


Con l'avvio dello Sprint 9, a partire dal 2026-06-02, il gruppo avvia
l'implementazione dimostrativa del Proof of Concept.

Il Proof of Concept ha lo scopo di dimostrare la fattibilità tecnica delle scelte
effettuate e la possibilità di soddisfare i requisiti principali individuati per
la Requirements \& Technology Baseline.

Il PoC\textsubscript{G} costituisce una baseline tecnologica e non architetturale: esso ha valore
dimostrativo e può essere rivisto nella successiva fase PB, qualora le
evidenze raccolte lo rendano opportuno.

La codifica svolta in questa fase non costituisce ancora il processo completo di
sviluppo software previsto per la fase PB.

\paragraph{Perimetro del PoC}

Il perimetro minimo del PoC comprende:

\begin{itemize}
    \item editor Markdown basato su CodeMirror;
    \item visualizzazione split tra editor e anteprima;
    \item rendering dell'anteprima Markdown;
    \item applicazione della formattazione in grassetto;
    \item salvataggio locale di contenuti Markdown;
    \item importazione di contenuti Markdown;
    \item apertura dell'editor su una nota esistente;
    \item interfaccia per l'interazione con LLM;
    \item popup di proposta generata dall'LLM;
    \item azioni di accettazione e rifiuto della proposta;
    \item selezione del cappello da applicare al testo selezionato;
    \item interfaccia di Distant Writing con inserimento del prompt\textsubscript{G}.
\end{itemize}

Il PoC non prevede, in fase RTB, persistenza su database. La gestione dei
contenuti è limitata a dati locali o file Markdown.

\paragraph{Tecnologie del PoC}

Il PoC utilizza:

\begin{itemize}
    \item Vue.js per il frontend;
    \item CodeMirror per l'editor Markdown;
    \item FastAPI per il backend;
    \item Docker Compose per l'esecuzione coordinata dei servizi;
    \item due container: \texttt{frontend} e \texttt{backend}.
\end{itemize}

La struttura del PoC è collocata nella cartella \texttt{poc/} della repository
\texttt{Second-Brain}, con separazione tra frontend e backend.

\paragraph{Criteri di successo}

Il PoC è considerato soddisfacente se consente di:

\begin{itemize}
    \item avviare localmente frontend e backend tramite Docker Compose;
    \item dimostrare la comunicazione minima tra frontend e backend;
    \item visualizzare un editor Markdown funzionante;
    \item mostrare l'anteprima del contenuto Markdown tramite visualizzazione split;
    \item applicare almeno una formattazione al testo;
    \item mostrare l'interfaccia prevista per l'interazione con LLM;
    \item mostrare il flusso di accettazione o rifiuto di una proposta generata;
    \item mostrare l'interfaccia di Distant Writing con inserimento del prompt;
    \item raccogliere evidenze utili alla progettazione di dettaglio in PB;
    \item individuare eventuali rischi tecnici da mitigare nella fase PB.
\end{itemize}

\paragraph{Validazione del PoC}

Il PoC viene validato internamente dal Progettista e dal Verificatore.

Il confronto con la proponente avviene prima della RTB ed è gestito dal
Responsabile, con il supporto dei membri che hanno lavorato sugli aspetti tecnici
e funzionali coinvolti.

I risultati del Proof of Concept devono:

\begin{itemize}
    \item essere discussi internamente;
    \item essere confrontati con la proponente;
    \item essere presentati durante la Requirements \& Technology Baseline;
    \item guidare la successiva progettazione di dettaglio in PB.
\end{itemize}

Eventuali osservazioni della proponente devono essere riportate in verbale esterno
e, se producono modifiche operative, trasformate in issue.

\paragraph{Sviluppo del Prodotto (Fase PB)}
Con l'ingresso nella fase PB, l'attività di codifica si estende allo sviluppo del prodotto finale (Minimum Viable Product). Per garantire uniformità e coerenza nello sviluppo collaborativo, il gruppo ha formalizzato regole ferree per la stesura del codice.

\paragraph{Perimetro del MVP}

A differenza del perimetro minimo del PoC, il Minimum Viable Product copre la
totalità dei requisiti obbligatori individuati nell'Analisi dei Requisiti,
tra cui:

\begin{itemize}
    \item editor Markdown completo (formattazione, elenchi, tabelle, link, undo/redo);
    \item riassunto, traduzione, riscrittura e Distant Writing tramite LLM;
    \item analisi critica secondo tutti e sei i ``Cappelli per Pensare'' di Edward de Bono;
    \item accettazione, rifiuto e rigenerazione della proposta AI, con interruzione manuale;
    \item salvataggio, apertura ed esportazione (PDF, HTML, JSON) della nota su file system locale, incluso il fallback per browser privi di File System Access API;
    \item gestione strutturata degli errori di comunicazione con i servizi esterni.
\end{itemize}

I requisiti opzionali relativi alla persistenza remota e alla gestione utenti
(R4-V-Op, R78--R86) restano fuori da questo perimetro (dettaglio nella
\doclink{https://synergo-unit-unipd.github.io/Second-Brain/docs/PB/doc_tecnici/doc_esterni/specifica_tecnica/specifica_tecnica.pdf}{Specifica Tecnica}, sezione Tracciamento e Stato dei Requisiti).

\paragraph{Tecnologie del MVP}

Il MVP utilizza lo stack consolidato durante la RTB ed esteso in PB:

\begin{itemize}
    \item Vue.js 3 e TypeScript per il frontend;
    \item CodeMirror per l'editor Markdown;
    \item FastAPI e Python per il backend;
    \item Docker Compose per l'esecuzione coordinata dei servizi;
    \item due container: \texttt{frontend} e \texttt{backend}.
\end{itemize}

Il dettaglio completo delle tecnologie e delle relative motivazioni è riportato
nella Specifica Tecnica. La struttura del MVP è collocata nella cartella
\texttt{mvp/} della repository \texttt{Second-Brain}, distinta dalla cartella
\texttt{poc/} che conserva il codice del Proof of Concept.

\paragraph{Criteri di successo del MVP}

Il MVP è considerato soddisfacente se:

\begin{itemize}
    \item soddisfa il 100\% dei requisiti obbligatori individuati nell'Analisi dei Requisiti;
    \item supera la totalità dei test unitari, di integrazione e di sistema pianificati (Sezione~\ref{sec:verifica});
    \item rispetta le soglie minime di code coverage definite (90\% Backend, 85\% Frontend);
    \item la pipeline di Continuous Integration risulta verde su ogni Pull Request integrata nel main;
    \item è accompagnato da Manuale Utente e Specifica Tecnica completi.
\end{itemize}

\paragraph{Validazione del MVP}

Il MVP viene validato internamente dal Verificatore, sulla base dell'esito della
pipeline di CI e della Mappatura Test di Sistema - Casi d'Uso (Sezione~\ref{sec:validazione}).
Il confronto con la proponente avviene durante la Product Baseline ed è gestito
dal Responsabile, con il supporto dei membri che hanno lavorato sugli aspetti
tecnici e funzionali coinvolti; eventuali osservazioni della proponente sono
riportate in verbale esterno e, se producono modifiche operative, trasformate
in issue, secondo lo stesso schema già adottato per la validazione del PoC.

\paragraph{Convenzioni di Nomenclatura (Naming)}
Tutti i membri, nel ruolo di Programmatore, devono attenersi rigorosamente alle seguenti convenzioni:
\begin{itemize}
    \item \textbf{Classi (TS e Python):} \texttt{PascalCase}.
    \item \textbf{Interfacce/Servizi (TS):} Nome del ruolo, senza prefisso "I" o suffissi tecnici (es. \texttt{NoteService}, non \texttt{INoteService}).
    \item \textbf{Adapter e Proxy:} Suffisso \texttt{Adapter} o \texttt{Proxy} in base al pattern (es. \texttt{OpenAIAdapter}, \texttt{NoteServiceProxy}).
    \item \textbf{Metodi e Variabili (TypeScript):} \texttt{camelCase}.
    \item \textbf{Moduli e Funzioni (Python):} \texttt{snake\_case} in accordo allo standard PEP8.
    \item \textbf{File di Test:} Devono avere lo stesso nome della classe testata, con suffisso \texttt{.test.ts} per il frontend e prefisso \texttt{test\_*.py} per il backend.
\end{itemize}

% ------------------------------------------------

\subsubsection{Integrazione del software}

\paragraph{Scopo}

Integrazione delle unità software nel "software item" completo e relativi test.

\paragraph{Workflow Git e CI/CD}
L'integrazione delle unità software avviene attraverso il modello di trunk-based development, nel quale le modifiche confluiscono nel branch \texttt{main} seguendo un workflow a passi fissi:
 
\begin{enumerate}
    \item lo sviluppo avviene su un branch dedicato, a partire dal quale i test vengono eseguiti in locale;
    \item una volta che i test locali sono passati, il branch viene pushato e si apre una Pull Request in draft, che avvia automaticamente la pipeline di CI;
    \item a termine dello sviluppo, viene compilata e spuntata la checklist della Pull Request, che viene quindi contrassegnata come pronta per la revisione (\textit{ready for review}), uscendo dallo stato di draft;
    \item solo a quetso punto avviene l'azione di verifica da parte del Verificatore, che approva la Pull Request unicamente se la pipeline risulta interamente verde.
\end{enumerate}
 
A partire dagli Sprint 12 e 13 è stato avviato il setup dell'ambiente di Continuous Integration/Continuous Deployment (CI/CD, GitHub Actions); il workflow sopra descritto è oggi pienamente operativo e obbligatorio per ogni Pull Request (Sezione~\ref{sec:verifica}).
\paragraph{Test Unitari (TU)}
Ogni unità software (classe o modulo) è accompagnata da test unitari automatici che ne verificano il comportamento in isolamento, con l'ausilio di doppi di test (mock, stub) per le dipendenze esterne. I test unitari sono eseguiti tramite \texttt{pytest} (Backend) e \texttt{vitest} (Frontend), integrati nella pipeline di Continuous Integration (Sezione~\ref{sec:verifica}).

\paragraph{Test di Integrazione (TI)}
L'efficacia dell'unione dei moduli è validata tramite appositi Test di Integrazione (TI) volti a controllare la comunicazione tra le unità (es. Frontend e Backend, Backend e servizio LLM esterno). A completamento della fase PB, la suite di Test di Integrazione è stata realizzata ed eseguita sistematicamente in pipeline, secondo la soglia di copertura definita nel Processo di Verifica (Sezione~\ref{sec:verifica}).

% ------------------------------------------------

\subsubsection{Test di qualifica del software}

\paragraph{Scopo}

Il test di qualifica del software verifica che il prodotto software soddisfi i propri requisiti.

\paragraph{Mappatura UC-Test}

La verifica dei requisiti software si basa sulla Mappatura Test di Sistema - Casi d’Uso, che definisce un collegamento univoco tra ogni funzionalità individuata nell'Analisi dei Requisiti e uno specifico test di sistema (TS). 

\paragraph{Metriche di Prodotto}

Il superamento di tali test deve essere accompagnato dal monitoraggio di metriche oggettive, tra cui il Soddisfacimento dei Requisiti Obbligatori (MPD01), il cui valore deve attestarsi tassativamente al 100\% per garantire la piena conformità del prodotto alle specifiche concordate con la proponente.

% ------------------------------------------------

\subsubsection{Integrazione del sistema}

\paragraph{Scopo}

Integrazione del software con hardware e altre componenti di sistema. 

\paragraph{Interazione con attori esterni}

Sebbene il progetto sia prettamente software, l'integrazione di sistema riguarda la connessione sinergica tra Frontend (Vue.js), Backend (FastAPI) e i servizi LLM esterni (Gemma, ChatGPT-OSS o Qwen) tramite API aderenti allo standard OpenAI.

\paragraph{Containerizzazione}

Tale processo è facilitato e validato attraverso l'uso di Docker Compose, che permette l'esecuzione coordinata dei diversi servizi in ambienti containerizzati isolati e riproducibili.

% ------------------------------------------------

\subsubsection{Test di qualifica del sistema}

\paragraph{Scopo}

Il test di qualifica del sistema è il test finale per dimostrare che il sistema completo è pronto per la consegna.

\paragraph{Validazione della Technology Baseline}

I test di qualifica del sistema in fase RTB sono finalizzati alla validazione della Technology Baseline. Essi dimostrano che l'intera infrastruttura (interfaccia utente, logica di backend e interazione AI) funzioni correttamente secondo gli scenari d'uso principali descritti nell'Analisi dei Requisiti.

\paragraph{Validazione della Product Baseline}

In fase PB, i test di qualifica del sistema sono stati eseguiti sul Minimum Viable Product completo, secondo la Mappatura Test di Sistema - Casi d'Uso (Sezione~\ref{sec:validazione}), dimostrando il soddisfacimento della totalità dei requisiti obbligatori individuati nell'Analisi dei Requisiti e coprendo, oltre agli scenari applicativi principali, anche i percorsi di errore e i casi limite (validazione dimensioni tabella, annullamento utente, indisponibilità del servizio LLM).

% ------------------------------------------------

\subsubsection{Installazione del software}

\paragraph{Scopo}

L'installazione del software è la pianificazione ed esecuzione dell'installazione nell'ambiente target.

\paragraph{Manuale Tecnico (README) e Manuale Utente}

La procedura di installazione e configurazione nell'ambiente target è documentata nel file README presente nella cartella \texttt{mvp/} del repository, ad uso di chi debba avviare o sviluppare il sistema. Tale documento riporta i prerequisiti necessari, la procedura di avvio tramite Docker Compose e le modalità di verifica del corretto funzionamento dell'intero sistema su host locale.

Ad uso dell'utente finale, in fase PB è stato inoltre redatto il Manuale Utente, che descrive in linguaggio non tecnico l'installazione, la configurazione e l'utilizzo di ciascuna funzionalità del prodotto, inclusa la risoluzione dei problemi più comuni.

% ------------------------------------------------

\subsubsection{Supporto all'accettazione del software}

\paragraph{Scopo}

Il supporto all'accettazione del software riguarda l'attività di supporto al cliente durante le revisioni e i test di accettazione finale.

\paragraph{Validazione con la Proponente}

Il processo prevede inoltre incontri di validazione periodici con la proponente Zucchetti, i cui esiti (accordi e correzioni richieste) sono formalizzati in appositi verbali esterni per guidare l'approvazione finale del prodotto.